\documentclass[10pt,a4paper]{amsart}
\usepackage{amsmath,amssymb,amsthm,graphicx}
\usepackage{verbatim}
\newenvironment{reference}{\comment}{\endcomment}
%\newenvironment{reference}{}{}
\newenvironment{slogan}{\comment}{\endcomment}
\newenvironment{history}{\comment}{\endcomment}

% For commutative diagrams we use Xy-pic
\usepackage[all]{xy}

% We use 2cell for 2-commutative diagrams.
\xyoption{2cell}
% Theorem environments.
%
\theoremstyle{plain}
\newtheorem{theorem}[subsection]{Theorem}
\newtheorem{proposition}[subsection]{Proposition}
\newtheorem{lemma}[subsection]{Lemma}

\theoremstyle{definition}
\newtheorem{definition}[subsection]{Definition}
\newtheorem{example}[subsection]{Example}
\newtheorem{exercise}[subsection]{Exercise}
\newtheorem{situation}[subsection]{Situation}

\theoremstyle{remark}
\newtheorem{remark}[subsection]{Remark}
\newtheorem{remarks}[subsection]{Remarks}

\numberwithin{equation}{subsection}

% Macros
%
\def\lim{\mathop{\mathrm{lim}}\nolimits}
\def\colim{\mathop{\mathrm{colim}}\nolimits}
\def\Spec{\mathop{\mathrm{Spec}}}
\def\Hom{\mathop{\mathrm{Hom}}\nolimits}
\def\Ext{\mathop{\mathrm{Ext}}\nolimits}
\def\SheafHom{\mathop{\mathcal{H}\!\mathit{om}}\nolimits}
\def\SheafExt{\mathop{\mathcal{E}\!\mathit{xt}}\nolimits}
\def\Sch{\mathit{Sch}}
\def\Mor{\mathop{\mathrm{Mor}}\nolimits}
\def\Ob{\mathop{\mathrm{Ob}}\nolimits}
\def\Sh{\mathop{\mathit{Sh}}\nolimits}
\def\NL{\mathop{N\!L}\nolimits}
\def\CH{\mathop{\mathrm{CH}}\nolimits}
\def\proetale{{pro\text{-}\acute{e}tale}}
\def\etale{{\acute{e}tale}}
\def\QCoh{\mathit{QCoh}}
\def\Ker{\mathop{\mathrm{Ker}}}
\def\Im{\mathop{\mathrm{Im}}}
\def\Coker{\mathop{\mathrm{Coker}}}
\def\Coim{\mathop{\mathrm{Coim}}}

% Boxtimes
%
\DeclareMathSymbol{\boxtimes}{\mathbin}{AMSa}{"02}

%
% Macros for moduli stacks/spaces
%
\def\QCohstack{\mathcal{QC}\!\mathit{oh}}
\def\Cohstack{\mathcal{C}\!\mathit{oh}}
\def\Spacesstack{\mathcal{S}\!\mathit{paces}}
\def\Quotfunctor{\mathrm{Quot}}
\def\Hilbfunctor{\mathrm{Hilb}}
\def\Curvesstack{\mathcal{C}\!\mathit{urves}}
\def\Polarizedstack{\mathcal{P}\!\mathit{olarized}}
\def\Complexesstack{\mathcal{C}\!\mathit{omplexes}}
% \Pic is the operator that assigns to X its picard group, usage \Pic(X)
% \Picardstack_{X/B} denotes the Picard stack of X over B
% \Picardfunctor_{X/B} denotes the Picard functor of X over B
\def\Pic{\mathop{\mathrm{Pic}}\nolimits}
\def\Picardstack{\mathcal{P}\!\mathit{ic}}
\def\Picardfunctor{\mathrm{Pic}}
\def\Deformationcategory{\mathcal{D}\!\mathit{ef}}
\usepackage{cleveref}
\crefname{originallemma}{original lemma}{original lemmas}
\Crefname{originallemma}{Original lemma}{Original lemmas}
\makeatletter
\newlabel{lemma-site-point-morphism}{{04EL}{source}{}{sourcepointer-1}{}}
\newlabel{lemma-site-point-morphism@cref}{{[originallemma][0][]04EL}{[1][0][]source}{}{}{}}
\makeatother
\crefname{originallemma}{original lemma}{original lemmas}
\Crefname{originallemma}{Original lemma}{Original lemmas}
\makeatletter
\newlabel{lemma-localize-morphism}{{03CF}{source}{}{sourcepointer-2}{}}
\newlabel{lemma-localize-morphism@cref}{{[originallemma][0][]03CF}{[1][0][]source}{}{}{}}
\makeatother
\crefname{originallemma}{original lemma}{original lemmas}
\Crefname{originallemma}{Original lemma}{Original lemmas}
\makeatletter
\newlabel{lemma-cocontinuous-morphism-topoi}{{00XO}{source}{}{sourcepointer-3}{}}
\newlabel{lemma-cocontinuous-morphism-topoi@cref}{{[originallemma][0][]00XO}{[1][0][]source}{}{}{}}
\makeatother
\crefname{originalsection}{original section}{original sections}
\Crefname{originalsection}{Original section}{Original sections}
\makeatletter
\newlabel{section-cocontinuous-adjoint}{{00XW}{source}{}{sourcepointer-4}{}}
\newlabel{section-cocontinuous-adjoint@cref}{{[originalsection][0][]00XW}{[1][0][]source}{}{}{}}
\makeatother
\title{The point of a localized site induced by a point and an element}
\author{The Stacks Project Authors}
\date{stacks-project:sites.tex:local-snapshot-20261012; GNU FDL 1.2 or later}
\begin{document}
\maketitle
\noindent\textit{Editorial note.} All delivered mathematical text is the author's own, copied from the pinned Stacks Project chapter file sites-cohomology.tex; no statement or proof step was rewritten, reordered or removed.\par\medskip
\noindent\textit{Editorial note.} The proof is the author's complete argument; every step the author writes out is inside the counted span and no closing sentence is a placeholder.\par\medskip
\noindent\textit{Editorial note.} Every labelled reference occurring inside the counted spans is carried in assembly.external\_references as a non-load-bearing navigation pointer (same-chapter ones at their definition site, cross-chapter ones at the invocation inside the counted span), so no load-bearing mathematics is replaced by a stub. No bibliography citation occurs inside the counted spans.\par\medskip
\noindent\textit{Editorial note.} The checked chapter list for same-name collisions is the set of stacks chapters staged in this repository: chow, cohomology, divisors, etale-cohomology, flat, more-algebra, more-morphisms, morphisms, spaces-more-morphisms, varieties together with this card's chapter sites-cohomology.\par\medskip
\noindent\textit{Editorial note.} The preamble is the chapter's own retained minimal preamble (sources/stacks-sites-cohomology/retained-preamble.tex), a verbatim excerpt of the native preamble: the Xy-pic setup and the native theorem declarations with the mathematical macros.\par\medskip
\noindent\textit{Allowed standard prerequisites (editorial).}\begin{itemize}
\item Points of a site as morphisms from the topos of a point. Same chapter, Lemma lemma-site-point-morphism. Represents the given point by a set-valued functor and identifies the topos of a point with the topos of the associated category of sets.
\item Localizing a morphism of sites. Same chapter, Lemma lemma-localize-morphism. Produces the localized morphism of sites fitting into the commutative diagram of topoi.
\end{itemize}\medskip
\ifdefined\hypertarget\hypertarget{sourcepointer-1}{}\fi
\noindent\textit{Source reference.} Original lemma 04EL: sources/stacks-sites/pinned-original-source.tex, source lines 8908--8908. Navigation pointer invoked at proof line 8908; the counted span contains the pointer only and the cited result is not reproduced inside it, so the pointer is registered as a non-load-bearing navigation stub in the established form for Stacks material. No bibliography citation is involved.\par\medskip
\ifdefined\hypertarget\hypertarget{sourcepointer-2}{}\fi
\noindent\textit{Source reference.} Original lemma 03CF: sources/stacks-sites/pinned-original-source.tex, source lines 8914--8914. Navigation pointer invoked at proof line 8914; the counted span contains the pointer only and the cited result is not reproduced inside it, so the pointer is registered as a non-load-bearing navigation stub in the established form for Stacks material. No bibliography citation is involved.\par\medskip
\ifdefined\hypertarget\hypertarget{sourcepointer-3}{}\fi
\noindent\textit{Source reference.} Original lemma 00XO: sources/stacks-sites/pinned-original-source.tex, source lines 8953--8953. Navigation pointer invoked at proof line 8953; the counted span contains the pointer only and the cited result is not reproduced inside it, so the pointer is registered as a non-load-bearing navigation stub in the established form for Stacks material. No bibliography citation is involved.\par\medskip
\ifdefined\hypertarget\hypertarget{sourcepointer-4}{}\fi
\noindent\textit{Source reference.} Original section 00XW: sources/stacks-sites/pinned-original-source.tex, source lines 8955--8955. Navigation pointer invoked at proof line 8955; the counted span contains the pointer only and the cited result is not reproduced inside it, so the pointer is registered as a non-load-bearing navigation stub in the established form for Stacks material. No bibliography citation is involved.\par\medskip
\begin{lemma}
\label{lemma-point-localize}
Let $\mathcal{C}$ be a site. Let $p$ be a point of $\mathcal{C}$ given by
$u : \mathcal{C} \to \textit{Sets}$. Let $U$ be an object of $\mathcal{C}$
and let $x \in u(U)$. The functor
$$
v : \mathcal{C}/U \longrightarrow \textit{Sets}, \quad
(\varphi : V \to U) \longmapsto \{y \in u(V) \mid u(\varphi)(y) = x\}
$$
defines a point $q$ of the site $\mathcal{C}/U$ such that the diagram
$$
\xymatrix{
& \Sh(pt) \ar[d]^p \ar[ld]_q \\
\Sh(\mathcal{C}/U) \ar[r]^{j_U} &
\Sh(\mathcal{C})
}
$$
commutes. In other words
$\mathcal{F}_p = (j_U^{-1}\mathcal{F})_q$ for any
sheaf on $\mathcal{C}$.
\end{lemma}

\begin{proof}
Choose $S$ and $\mathcal{S}$ as in
Lemma \ref{lemma-site-point-morphism}.
We may identify $\Sh(pt) = \Sh(\mathcal{S})$
as in that lemma, and we may write
$p = f : \Sh(\mathcal{S}) \to \Sh(\mathcal{C})$
for the morphism of topoi induced by $u$.
By
Lemma \ref{lemma-localize-morphism}
we get a commutative diagram of topoi
$$
\xymatrix{
\Sh(\mathcal{S}/u(U)) \ar[r]_-{j_{u(U)}} \ar[d]_{p'} &
\Sh(\mathcal{S}) \ar[d]^p \\
\Sh(\mathcal{C}/U) \ar[r]^{j_U} &
\Sh(\mathcal{C}),
}
$$
where $p'$ is given by the functor $u' : \mathcal{C}/U \to \mathcal{S}/u(U)$,
$V/U \mapsto u(V)/u(U)$.
Consider the functor $j_x : \mathcal{S} \cong \mathcal{S}/x$ obtained
by assigning to a set $E$ the set $E$ endowed with the constant map
$E \to u(U)$ with value $x$.
Then $j_x$ is a fully faithful cocontinuous functor which has a
continuous right adjoint
$v_x : (\psi : E \to u(U)) \mapsto \psi^{-1}(\{x\})$.
Note that $j_{u(U)} \circ j_x = \text{id}_\mathcal{S}$, and
$v_x \circ u' = v$.
These observations imply that we have the following commutative
diagram of topoi
$$
\xymatrix{
\Sh(\mathcal{S}) \ar[rd]^a \ar[rdd]_q
\ar `r[rrr] `d[dd]^p [rrdd] & & & \\
& \Sh(\mathcal{S}/u(U)) \ar[r]_-{j_{u(U)}} \ar[d]^{p'} &
\Sh(\mathcal{S}) \ar[d]^p & \\
& \Sh(\mathcal{C}/U) \ar[r]^{j_U} &
\Sh(\mathcal{C}) &
}
$$
Namely:
\begin{enumerate}
\item The morphism
$a : \Sh(\mathcal{S}) \to \Sh(\mathcal{S}/u(U))$
is the morphism of topoi associated to the cocontinuous functor
$j_x$, which equals the morphism associated to the continuous
functor $v_x$, see
Lemma \ref{lemma-cocontinuous-morphism-topoi}
and
Section \ref{section-cocontinuous-adjoint}.
\item The composition $p \circ j_{u(U)} \circ a = p$ since
$j_{u(U)} \circ j_x = \text{id}_\mathcal{S}$.
\item The composition $p' \circ a$ gives a morphism of topoi.
Moreover, it is the morphism of topoi associated to the continuous
functor $v_x \circ u' = v$. Hence $v$ does indeed define a point $q$ of
$\mathcal{C}/U$ which fits into the diagram above by construction.
\end{enumerate}
This ends the proof of the lemma.
\end{proof}

\end{document}
